Having highlighted in Chapter \ref{chap:discussion} what this thesis contributes to theory, practice, and methodology, I cannot help but feel a little strange writing about such contributions as if they belonged in separate boxes. This reflection is not meant to be generalizable, and it would certainly fail an external validity test. Then again, it’s meant as reflection, not replication. I reflect that none of these insights came from pursuing any contribution in isolation. Over five years of sweat and tears, I realized that they came from a messy, nonlinear process. Sometimes they came accidentally, sometimes intentionally, but most of the time procedurally. This is because the process by which I arrived at the contributions was a process made up of cycles. The cycles involved understanding the data, making sense of them, recognizing limitations, building visual tools, refining models, and then interpreting the results and the story they seemed to tell.

To make sense of anything, I first had to understand the data—not just statistically, but structurally and contextually. What did each variable represent? What patterns were already there before I started drawing lines through the data or asking algorithms to discover them for me? What was missing, why, and what could be done about it? These questions became the starting point for every project. And just as much as theory guides analysis, the data often forced me to revisit assumptions and rethink the questions I was asking. Eventually, I learned to ask better questions, shifting from “Is this effect significant?” to “What does this tell us about how people behave (or don’t) in certain situations, and about how businesses respond?”

Throughout my studies, some of the most striking insights did not come from rigorous models or perfect specifications. And by striking, I don’t mean the most statistically robust. Those came later, after the models were built and the story had already been laid out. The real insights came while exploring data through interactive dashboards that I designed myself. They allowed me to visualize patterns, compare segments, and trace how key variables evolved over time through colors, shapes, and simple interactive elements such as filters and dropdowns. No hypothesis tests, no coefficients; just time series, simple cohort comparisons, and a thousand tabs open at once. 

These tools were not merely exploratory. They suggestively pointed toward observed patterns. They didn’t prove the patterns, and working with them was less about analyzing the data, more about narrating it, about trying to see where meaning might be hiding. It felt like holding a book of theory in the dark, with the dashboard offering just enough light to read a few lines at a time. Some might say it’s enough to read in the dark, if you just focus hard enough. Others might say it’s all about finding the light and seeing what appears. But I think I prefer a little light to guide the reading.

Looking back on this process, I have come to realize that research in the social sciences is not just about technical proficiency or theoretical advancement. It is about learning to navigate contradictions between theory and data, detail and direction, explanation and uncertainty. It is about knowing when to zoom in on a variable and when to step back and ask if that variable matters at all. The research I conducted became less a process of control and more a negotiation with the data, the related literature, and my own assumptions. And while some might prefer to treat assumptions as something objective or at least external to the researcher, I’ve learned that they are anything but. They shape what we see, what we ignore, and what we are willing to believe the data is telling us.

Part of the reason I came to see research this way is that this thesis took long enough for me to spend some time in industry. The more I moved between the two (full time at the bank and part time at the PhD office), the stranger it felt to hear people talk about academia versus industry as if they constituted different worlds. From my experience, the academic process is almost identical to the industry process. Both try to make sense of observed behavior, to predict outcomes and improve decisions, whether through behavioral explanations or algorithmic models. The only real difference is emphasis. Academia polishes its explanations until they are publishable. Industry polishes its tools until they are profitable. One asks what is going on and why. The other asks what to do next. Both chase relevance, each through a different lens: the microscope of academia and the binoculars of industry.

That said, I can’t resist a confession.

Sometimes I wonder if social science has (or will) become a full-blown circus of triviality. Everyone’s performing, juggling regressions, walking the tightrope of robustness checks, and pulling contributions out of hats. We spend years polishing results so safe, so painfully obvious, that any trace of something even slightly interesting is lost under layers of statistical rigor. Then comes the best part: we frame the findings in terms of what they add to the literature, even though it always remains the same literature, in the same journals, read by the same people. As if the point of research were not to understand the world but to impress insiders. The only real novelty? A new dataset, usually slightly bigger or slightly weirder than the last one. It is exhausting. We are rewarded for saying something new (not necessarily something true, and not always something useful). At times, I feel like we are just optimizing for clever footnotes on the margins of irrelevance.

But let’s not pretend industry isn’t running its own sideshow. If academia is a circus of triviality, then industry is a magic show of illusory certainty, of quick solutions pulled out of black box algorithms, of KPIs waved like proof of wisdom, of ambiguity made to disappear behind convincing tables and dashboards. Industry frames impact but fears uncertainty, celebrates innovation but avoids reflection. The costumes differ, but the performance is the same: a pursuit of efficiency over understanding, staged according to the same kind of instrumentalism.

And yet, I remain optimistic. Because the real value, I think, lies in the overlap. In realizing that research is not just about what practitioners can learn from us, but what we can learn from them. People in industry do not need to be told context matters; they live in it. They do not theorize about heterogeneity; they market to it. And maybe, just maybe, good research is not the kind that holds itself apart from practice. Maybe it’s the kind that builds practice from tools, constraints, and theory, lit by a dashboard like the one that sparked every idea in this thesis.
